\documentclass[../main.tex]{subfiles}
\begin{document}

\section{Numerical Validation}
\label{sec:validation}

Every closed form in Section~\ref{sec:theory} is checked against Monte-Carlo
simulation of Definition~\ref{def:ccsbm}. Our purpose here is verification of the
analysis, not a performance claim: the results below establish that the derivations
are correct, and say nothing on their own about accuracy on real bot graphs.
Section~\ref{sec:limitations} addresses that gap directly.

\paragraph{Protocol.}
Unless stated otherwise we use $\pi_B = 0.3$ ($c \approx 0.4286$),
$\|\Delta\| = 2$, $\sigma = 1$, and $D = 16$, drawing $4\times 10^5$ nodes per class
per configuration ($1.5 \times 10^6$ for the crossover sweep). Errors are computed
by exhaustive threshold search over the empirical distributions, so the reported
figures are attained minima rather than plug-in estimates. All code is released
with the paper and runs in a few minutes on CPU.

\subsection{Deflation and error}

Table~\ref{tab:deflation} compares Proposition~\ref{prop:deflation} and
Theorem~\ref{thm:error} against simulation. Separation matches to within
$1.2\times 10^{-3}$ and the standard deviation---including the mixture term that a
naive calculation omits---to within $4\times 10^{-4}$ across the whole range.
Figure~\ref{fig:deflation} shows the same agreement graphically.

\begin{table}[t]
\centering
\caption{Theory versus simulation for signal deflation and Bayes error
($\pi_B = 0.3$, $\|\Delta\|=2$, $\sigma=1$, $D=16$).}
\label{tab:deflation}
\small
\resizebox{\columnwidth}{!}{%
\begin{tabular}{@{}ccccccc@{}}
\toprule
& & \multicolumn{2}{c}{separation $\delta$} & \multicolumn{2}{c}{$s_{\mathrm{B}}$} & error \\
\cmidrule(lr){3-4}\cmidrule(lr){5-6}
$\rho$ & $\eta$ & theory & sim. & theory & sim. & thy / sim. \\
\midrule
0.00 & 0.000 & 2.0000 & 2.0000 & 0.2500 & 0.2499 & 0.00003 / 0.00003 \\
0.10 & 0.043 & 1.7143 & 1.7142 & 0.2915 & 0.2918 & 0.00113 / 0.00160 \\
0.20 & 0.086 & 1.4286 & 1.4283 & 0.3202 & 0.3204 & 0.00925 / 0.01095 \\
0.30 & 0.129 & 1.1429 & 1.1441 & 0.3391 & 0.3391 & 0.03698 / 0.03895 \\
0.40 & 0.171 & 0.8571 & 0.8586 & 0.3500 & 0.3499 & 0.09785 / 0.09898 \\
0.50 & 0.214 & 0.5714 & 0.5696 & 0.3536 & 0.3528 & 0.19898 / 0.19974 \\
0.60 & 0.257 & 0.2857 & 0.2864 & 0.3500 & 0.3500 & 0.33734 / 0.33680 \\
\bottomrule
\end{tabular}%
}
\end{table}

\begin{figure}[t]
    \centering
    \includegraphics[width=\linewidth]{figures/theory_deflation.png}
    \caption{Theory versus simulation. (a) Post-aggregation separation decays
    linearly in $\rho$ exactly as Proposition~\ref{prop:deflation} predicts.
    (b) Bayes error against the graph-free baseline; the predicted critical ratio
    $\rho^\star = 0.464$ coincides with the crossing point.}
    \label{fig:deflation}
\end{figure}

\subsection{The critical ratio}

The central prediction is Corollary~\ref{cor:critical}. Solving
$\mathrm{SNR}(\rho) = \|\Delta\|/(2\sigma)$ numerically at $D=16$ gives
$\rho^\star = 0.4644$. Sweeping $\rho$ in steps of $0.01$ with $1.5\times 10^6$
samples per point, aggregation beats raw features at $\rho = 0.46$
($\varepsilon = 0.1545$ against $\varepsilon_{\mathrm{raw}} = 0.1587$) and loses at
$\rho = 0.47$ ($\varepsilon = 0.1648$). The measured crossover therefore lies in
$[0.46, 0.47]$, containing the predicted value.

Table~\ref{tab:critical} separates the two forms of the corollary. The exact
crossover differs from the closed form~\eqref{eq:critical} at moderate SNR, as
expected, since~\eqref{eq:critical} drops the mixture variance. Re-running the
exact solver in the noise-dominated regime it was derived for ($\|\Delta\| = 0.2$,
$\sigma = 1$) recovers the closed form to three decimal places at every degree,
confirming that the discrepancy is the stated approximation and not an error.

\begin{table}[t]
\centering
\caption{Critical camouflage ratio: closed form~\eqref{eq:critical} against the
exact SNR crossover, at moderate and low SNR.}
\label{tab:critical}
\small
\begin{tabular}{@{}cccc@{}}
\toprule
$D$ & closed form~\eqref{eq:critical} & exact ($\|\Delta\|/\sigma = 2$) & exact (low SNR) \\
\midrule
4  & 0.3500 & 0.2604 & 0.3488 \\
9  & 0.4667 & 0.3913 & 0.4657 \\
16 & 0.5250 & 0.4644 & 0.5243 \\
25 & 0.5600 & 0.5102 & 0.5594 \\
64 & 0.6125 & 0.5807 & 0.6121 \\
\bottomrule
\end{tabular}
\end{table}

Figure~\ref{fig:critical} plots both curves against degree. The tolerated
camouflage rises steeply for sparsely connected accounts and flattens as $D$ grows,
which has a concrete operational reading: the accounts most vulnerable to
camouflage are those with few connections, and a deployment that aggregates over
small neighbourhoods is fragile in exactly the regime where bots are cheapest to
create.

\begin{figure}[t]
    \centering
    \includegraphics[width=0.72\linewidth]{figures/theory_critical_ratio.png}
    \caption{Critical camouflage ratio against neighbourhood size. Below the curve
    the graph improves on raw features; above it the graph is actively harmful.
    The closed form is the noise-dominated limit of the exact crossover.}
    \label{fig:critical}
\end{figure}

\subsection{SGTR}

Lemma~\ref{lem:energy} predicts expected squared discrepancies of $2d\sigma^2$ and
$\|\Delta\|^2 + 2d\sigma^2$. At $d = 8$, $\sigma = 1$, $\|\Delta\| = 2$ these are
$16$ and $20$; simulation over $2\times10^5$ edge samples gives $16.005$ and
$20.005$, an error below $0.03\%$.

For Theorem~\ref{thm:sgtr} we verify the closed
form~\eqref{eq:rhostarsgtr} by round-trip: substituting
$\rho^\star_{\mathrm{SGTR}}$ back into $\rho_{\mathrm{eff}}$ must return
$\rho^\star$. The maximum absolute deviation over
$\bar e_c \in \{0.25, 0.5, 0.75, 0.9\}$ is $1.1\times 10^{-16}$, i.e.\ machine
precision. Table~\ref{tab:sgtr} reports the resulting tolerance gains, and
Figure~\ref{fig:sgtr} plots the full curve.

\begin{table}[t]
\centering
\caption{Effect of SGTR attenuation on tolerated camouflage
($D=16$, $\pi_B=0.3$, baseline $\rho^\star = 0.525$).}
\label{tab:sgtr}
\small
\begin{tabular}{@{}ccc@{}}
\toprule
$\bar e_c$ & $\rho_{\mathrm{eff}}$ at $\rho = 0.4$ & $\rho^\star_{\mathrm{SGTR}}$ \\
\midrule
0.00 & 0.4000 & 0.5250 \\
0.25 & 0.3333 & 0.5957 \\
0.50 & 0.2500 & 0.6885 \\
0.75 & 0.1429 & 0.8155 \\
0.90 & 0.0625 & 0.9170 \\
\bottomrule
\end{tabular}
\end{table}

\begin{figure}[t]
    \centering
    \includegraphics[width=0.72\linewidth]{figures/theory_sgtr_tolerance.png}
    \caption{Tolerated camouflage as a function of the mean attenuation applied to
    camouflage edges. The shaded region is the operating range that spectral
    pruning adds over an unpruned homophilic aggregator.}
    \label{fig:sgtr}
\end{figure}

\subsection{The extensions: achievable pruning and depth}
\label{sec:validation-ext}

Proposition~\ref{prop:deff} and Theorem~\ref{thm:depth} are checked in
\texttt{theory/extensions2.py}. The effective-degree formula~\eqref{eq:deff} matches
Monte-Carlo variance to within $0.5\%$ for $\theta \geq 0.5$, degrading to $7\%$ at
$\theta = 0$ where finitely many nodes lose their entire neighbourhood and are
excluded. The perfect-pruning ceiling of Corollary~\ref{cor:ceiling} reproduces
$1-1/D$ to six decimal places for $D \in \{4,16,64,256\}$. The tolerance curve is
monotone in the retained-weight ratio, rising from $0.525$ at $\theta = 1$ to
$0.9375$ at $\theta = 0$ for $D=16$, $\pi_B = 0.3$.

For the depth recursion, the exact recursion
\eqref{eq:depth-mean}--\eqref{eq:depth-varH} agrees with a locally tree-like
simulation to between $0.1\%$ and $1.5\%$ across $\rho \in \{0,0.2,0.35,0.5\}$ and
$L \in \{1,2,3\}$. The simplified form~\eqref{eq:depth-clean} is exact at $\rho = 0$
but overestimates once the mixture terms matter---by a factor of two at $\rho=0.2$,
$L=3$---confirming that it is genuinely a noise-dominated limit and not the general
case; raising $\sigma$ to $6$ recovers it to within $4\%$. The two crossover points
of the exact recursion sit at $0.5245$ and $0.8754$ for $L = 1$, $0.5245$ and
$0.8755$ for $L=2$, and $0.5245$ and $0.8755$ for $L=3$, against the predicted
window $(0.525, 0.875)$: the crossover is depth-independent to four decimal places.

\subsection{Measuring the camouflage ratio on real graphs}
\label{sec:measuring}

The theory is parameterised by $\rho$, which has so far been a modelling device. We
now measure it. TwiBot-20 and TwiBot-22 require author approval for access, which we
did not have; we therefore use the two graphs from
CARE-GNN~\cite{dou2020caregnn}---the work that introduced the camouflage framing---
together with the Elliptic bitcoin graph, all openly available. These are fraud
rather than bot graphs. The mechanism is the same (a fraudulent node attaching to
benign ones to dilute its neighbourhood) but the domain is not, and we draw no
conclusion about Twitter from them.

For each relation we compute $\rho$ as the mean over fraud nodes of the fraction of
neighbours that are benign, and $\eta$ symmetrically. Table~\ref{tab:real} reports
the result.

\begin{table}[t]
\centering
\small
\caption{Camouflage measured on real labelled graphs. $\hat\rho$ is the mean benign
share of a fraud node's neighbourhood; the window is
$\bigl((1-D^{-1/2})/(1+c),\,(1+D^{-1/2})/(1+c)\bigr)$ from~\eqref{eq:window}
evaluated at the median fraud degree. Relations with median degree $1$ are omitted as
degenerate: with a single neighbour there is no variance reduction to be had.}
\label{tab:real}
\begin{tabular}{llrrrr}
\toprule
Graph & Relation & $\hat\rho$ & $\hat\eta$ & $c\hat\rho$ & Window \\
\midrule
YelpChi & \texttt{homo}    & 0.805 & 0.133 & 0.137 & (0.79, 0.92) \\
        & \texttt{net\_rtr} & 0.820 & 0.140 & 0.139 & (0.68, 1.03) \\
        & \texttt{net\_rsr} & 0.795 & 0.135 & 0.135 & (0.78, 0.93) \\
\midrule
Amazon  & \texttt{homo}    & 0.897 & 0.033 & 0.066 & (0.85, 1.01) \\
        & \texttt{net\_usu} & 0.897 & 0.036 & 0.066 & (0.84, 1.03) \\
        & \texttt{net\_uvu} & 0.818 & 0.030 & 0.060 & (0.67, 1.19) \\
\bottomrule
\end{tabular}
\end{table}

Three observations, and one caveat we consider important.

\paragraph{The edge-balance identity holds where its assumptions do.}
On YelpChi the prediction $\eta = c\rho$ is accurate to within $3\%$ on
\texttt{homo} ($0.137$ predicted against $0.133$ observed) and to three decimal
places on \texttt{net\_rsr} ($0.135$ against $0.135$). On Amazon it overestimates by
roughly a factor of two. The discrepancy is explained by degree heterogeneity:
\eqref{eq:balance} counts edge endpoints under the model's regular-degree assumption,
and generalises to $\eta = c\rho\,(\bar d_B/\bar d_H)$ when the classes have
different mean degree. YelpChi has $\bar d_B/\bar d_H \approx 0.95$--$1.01$, Amazon
$0.31$--$0.35$ on its dense relations. The corrected form improves Amazon's
\texttt{homo} prediction from $0.066$ to $0.023$ against an observed $0.033$, but
overshoots on \texttt{net\_upu} ($\bar d_B/\bar d_H = 1.80$), because $\eta$ is a
mean of per-node ratios rather than a ratio of means and the two differ under
within-class degree skew. We therefore report the identity as validated on graphs
with comparable class degrees and as requiring the degree correction otherwise.

\paragraph{Real fraud graphs sit inside the harmful window, marginally.}
Every dense relation in Table~\ref{tab:real} has $\hat\rho$ above the lower edge of
its window, so the theory predicts that \emph{plain mean aggregation} on these graphs
is worse than ignoring the topology. The margins are small---$0.805$ against a
threshold of $0.787$ on YelpChi \texttt{homo}---which is the interesting part. These
graphs sit close to the boundary, which is precisely the regime in which the choice
of aggregator decides the outcome, and it is consistent with the empirical fact that
camouflage-resistant architectures were needed on exactly these datasets.

\paragraph{Weak pruning suffices, and the theory says why.}
Using the measured feature separations ($r = 1.17$ over $d=32$ for YelpChi,
$r = 2.04$ over $d=25$ for Amazon), \eqref{eq:kappa} gives $\kappa = 0.086$ and
$0.295$, hence best achievable retained-weight ratios $\theta_{\min} = 0.77$ and
$0.43$. These are poor pruners. Yet substituting them
into~\eqref{eq:pruned-condition} raises the tolerance from $0.787$ to $0.827$ on
YelpChi \texttt{homo} and from $0.854$ to $0.930$ on Amazon \texttt{homo}---in both
cases past the measured $\hat\rho$. The prediction is therefore that even a weak
discrepancy-based filter flips these graphs from harmful to helpful, which is a
concrete, falsifiable account of why neighbour selection is the component that
matters on these benchmarks. It also explains why a \emph{learned} similarity is
worth the trouble: at $\kappa < 0.3$ the raw score is close to uninformative.

\paragraph{Caveat.}
The window is derived for a CSBM with regular degrees, isotropic Gaussian features,
and a single relation. Real graphs violate all three: degrees are heavy-tailed,
features are neither Gaussian nor isotropic, and the relations interact. The
numbers above should be read as locating these graphs approximately, not as a
certified verdict, and the agreement of $\eta$ with $c\rho$ on YelpChi is the one
quantity here that constitutes a genuine test of the theory rather than an
application of it. Measuring $\rho$ on TwiBot remains the natural next step and
requires only data access; \texttt{theory/measure\_rho.py} implements the estimator
and validates it against planted ground truth, recovering $\rho$ exactly and
confirming $\eta = c\rho$ on synthetic graphs with regular degrees.

\subsection{THCA}

Theorem~\ref{thm:thca} is constructive and we verify both halves programmatically.
The neighbourhoods $\{0,2\}$ and $\{1,1\}$ have identical means ($1.0$ and $1.0$)
and distinct maxima ($2.0$ and $1.0$), so mean aggregation cannot separate them
while the max channel can; and setting $\alpha = (1,0,0)$ reproduces the homophilic
channel exactly. Both checks pass.

\end{document}
